% --- Estilo de gráficas de datos (como en «Leyes básicas de la electrostática»)
\pgfplotsset{
  datos/.style={
    tick label style={font=\scriptsize}, label style={font=\small},
    axis lines=left, axis line style={grispizarra},
    grid=major, grid style={grisclaro, line width=0.4pt},
    every tick/.style={grispizarra},
    legend style={font=\scriptsize, draw=none, fill=white, fill opacity=0.9,
                  text opacity=1, cells={anchor=west}},
    legend cell align=left,
    /pgf/number format/use comma, /pgf/number format/1000 sep={\,},
    scaled ticks=false,
  },
  puntoA/.style={only marks, mark=*, mark size=2.1pt, azulprof,
    mark options={fill=azulprof, draw=white, line width=0.3pt},
    error bars/.cd, x dir=both, x explicit, y dir=both, y explicit,
    error bar style={line width=0.6pt, azulprof}, error mark options={rotate=90, mark size=2.2pt, line width=0.6pt, azulprof}},
  puntoB/.style={only marks, mark=square*, mark size=1.9pt, azulacero,
    mark options={fill=azulclaro, draw=azulacero, line width=0.5pt},
    error bars/.cd, x dir=both, x explicit, y dir=both, y explicit,
    error bar style={line width=0.6pt, azulacero}, error mark options={rotate=90, mark size=2.2pt, line width=0.6pt, azulacero}},
  ajuste/.style={azulprof, line width=1.1pt},
  teorica/.style={acento, line width=0.9pt, densely dashed},
}

% Referencia a un panel de figura: «figura 5\,(b)», enlazada en bloque
\newcommand{\figpanel}[2]{\hyperref[#1]{figura~\ref*{#1}\,\textup{(#2)}}}

